% Standalone problem document, generated from the adjacent README.md.
% Compile with XeLaTeX. Mathematical content is maintained in Markdown.
\documentclass[11pt,a4paper]{article}
\usepackage[fontsize=10.5pt]{scrextend}
\usepackage[margin=22mm,headheight=15pt,headsep=9mm,footskip=11mm]{geometry}
\usepackage{fontspec}
\setmainfont{texgyrepagella}[Extension=.otf,UprightFont=*-regular,BoldFont=*-bold,ItalicFont=*-italic,BoldItalicFont=*-bolditalic]
\setsansfont{texgyreheros}[Extension=.otf,UprightFont=*-regular,BoldFont=*-bold,ItalicFont=*-italic,BoldItalicFont=*-bolditalic]
\setmonofont{texgyrecursor}[Extension=.otf,UprightFont=*-regular,BoldFont=*-bold,ItalicFont=*-italic,BoldItalicFont=*-bolditalic,Scale=MatchLowercase]
\usepackage{amsmath,amssymb}
\usepackage{unicode-math}
\setmathfont{texgyrepagella-math.otf}
\usepackage{xcolor,microtype,enumitem,needspace,fancyhdr}
\usepackage{bookmark}
\definecolor{ink}{HTML}{17344B}
\definecolor{accent}{HTML}{197D80}
\definecolor{muted}{HTML}{596773}
\hypersetup{colorlinks=true,linkcolor=accent,urlcolor=accent,pdftitle={FR-03 - Uniform
conditioning of large Paley-frame column
subsets},pdfauthor={Open Problems in Numerical Linear Algebra}}
\urlstyle{same}
\setlength{\parindent}{0pt}
\setlength{\parskip}{5pt plus 1pt minus 1pt}
\linespread{1.04}
\setlength{\emergencystretch}{3em}
\widowpenalty=10000
\clubpenalty=10000
\displaywidowpenalty=10000
\allowdisplaybreaks[1]
\setlist{nosep,leftmargin=1.5em}
\providecommand{\tightlist}{\setlength{\itemsep}{0pt}\setlength{\parskip}{0pt}}
\providecommand{\pandocbounded}[1]{#1}
\pagestyle{fancy}
\fancyhf{}
\fancyhead[L]{\sffamily\footnotesize\color{muted}OPEN PROBLEMS IN NUMERICAL LINEAR ALGEBRA}
\fancyhead[R]{\sffamily\bfseries\footnotesize\color{accent}FR-03}
\fancyfoot[L]{\parbox[b]{0.9\textwidth}{\sffamily\fontsize{8}{10}\selectfont\color{muted}Editorial ratings; a bounded search cannot certify that a problem remains unsolved.\\Literature check: 2026-09-10}}
\fancyfoot[R]{\sffamily\footnotesize\color{muted}\thepage}
\renewcommand{\headrulewidth}{0.3pt}
\renewcommand{\headrule}{\hbox to\headwidth{\color{accent}\leaders\hrule height \headrulewidth\hfill}}
\makeatletter
\renewcommand{\section}{\Needspace{5\baselineskip}\@startsection{section}{1}{0pt}{12pt plus 2pt minus 2pt}{4pt}{\sffamily\bfseries\large\color{ink}}}
\renewcommand{\subsection}{\Needspace{4\baselineskip}\@startsection{subsection}{2}{0pt}{9pt plus 1pt minus 1pt}{3pt}{\sffamily\bfseries\normalsize\color{ink}}}
\makeatother
\setcounter{secnumdepth}{0}
\begin{document}
{\sffamily\small\color{muted}Frames and matrix designs\par}
\vspace{3pt}
{\raggedright\hyphenpenalty=10000\exhyphenpenalty=10000\sffamily\bfseries\fontsize{21}{24}\selectfont\color{ink}Uniform
conditioning of large Paley-frame column subsets\par}
\vspace{7pt}
{\sffamily\small\textbf{Difficulty:} extreme\quad\textbf{Importance:} broadly
interesting\par}
{\sffamily\small\bfseries\color{accent}Status: Open\par}
\vspace{4pt}
{\color{accent}\hrule height 0.8pt}
\vspace{6pt}
\textbf{Rating rationale:} Uniform conditioning for this explicit
arithmetic frame beyond square-root sparsity is a major deterministic
sensing barrier linking number theory and numerical reconstruction.

For a prime \(p\equiv1\pmod4\), let \(Q\) be the nonzero quadratic
residues modulo \(p\), and set \(N=(p+1)/2\). Define
\(\Phi_p\in\mathbb C^{N\times(p+1)}\) with row labels \(\{0\}\cup Q\)
and column labels \(\{0,\ldots,p-1,\infty\}\) by \[
(\Phi_p)_{0j}=p^{-1/2},\qquad
(\Phi_p)_{rj}=\sqrt{2/p}\exp(-2\pi\mathrm i rj/p)
\quad(r\in Q,\ 0\leq j<p),
\] and \((\Phi_p)_{0,\infty}=1\), \((\Phi_p)_{r,\infty}=0\) for
\(r\in Q\). All columns have Euclidean norm one.

Do constants \(0<\varepsilon<1/2\), \(c>0\), \(C<\infty\), and \(p_0\)
exist such that, for every prime \(p\geq p_0\) with \(p\equiv1\pmod4\)
and every column subset \(S\) with \[
1\leq |S|\leq \lfloor c p^{1/2+\varepsilon}\rfloor,
\] the column submatrix satisfies \[
\kappa_2((\Phi_p)_S)
=\frac{\sigma_{\max}((\Phi_p)_S)}
{\sigma_{\min}((\Phi_p)_S)}\leq C?
\] The condition number is infinite when the columns are dependent.
Constants must be independent of both \(p\) and \(S\).

This is the bounded-condition-number formulation of the Paley-frame
conjecture in Randomstrasse 101, Conjecture 29. It asks for uniform
conditioning beyond the square-root sparsity scale, rather than
conditioning of a random subset. The explicit frame would provide
deterministic measurement matrices with stable sparse least-squares
subproblems. The general deterministic restricted-isometry construction
problem does not require this particular matrix, so neither problem
duplicates the other.

\subsection{References}\label{references}

\begin{enumerate}
\def\labelenumi{\arabic{enumi}.}
\tightlist
\item
  A. S. Bandeira, D. Dmitriev, K. Lucca, P. Nizić-Nikolac, and A.
  Rödder, \emph{Randomstrasse 101: Open Problems of 2025},
  arXiv:2603.29571 (2026), Conjecture 29 and the preceding Paley
  construction. \href{https://arxiv.org/html/2603.29571v1}{Paper}.
\item
  A. S. Bandeira, M. Fickus, D. G. Mixon, and P. Wong, \emph{The road to
  deterministic matrices with the restricted isometry property}, Journal
  of Fourier Analysis and Applications 19 (2013), pp.~1123--1149. The
  Paley ETF discussion supplies the original restricted-isometry
  motivation. \href{https://arxiv.org/abs/1202.1234}{Paper}.
\item
  S. Satake, \emph{On the restricted isometry property of the Paley
  matrix}, Linear Algebra and its Applications 631 (2021), pp.~35--47.
  Its main implication is conditional on a Paley-graph conjecture.
  \href{https://arxiv.org/abs/2011.02907}{Paper}.
\end{enumerate}

\subsection{Status check --- 2026-09-10}\label{status-check-2026-09-10}

Searched ``Paley matrix RIP conjecture 2025 2026'', ``Paley ETF
condition number square root bottleneck'', and the exact Satake title;
checked the 2026 restatement and current arXiv records. No unconditional
proof or counterexample was found. Conditional graph-discrepancy
implications do not settle the quantifiers above.

\textbf{Audit update (2026-09-10):} Rechecked the 2026 Conjecture 29 and
searched for later Paley ETF results. The exact uniform condition-number
target remains explicitly conjectured; no unconditional resolution was
located. This is a bounded literature check, not a proof that no
solution exists.
\end{document}
